\subsection{Coxeter 群的子群}

在本号中，我们假定 \((W, S)\) 是一个 Coxeter 系统。对于 S 的任意子集 X，记 \(W_{X}\) 为由 X 生成的 W 的子群。

命题 7. 设 \(w\) 属于 W。存在一个子集 \(S_w\)，它是 \(S\) 的子集，使得 \(\{s_1, \ldots, s_q\} = S_w\) 对任意约化分解 \((s_1, \ldots, s_q)\) 的 \(w\) 都成立。

令 M 为由 S 的子集组成的幺半群，其运算为 \((A,B)\mapsto A\cup B\)；M 的单位元为 \(\varnothing\)。置 \(f(s)=\{s\}\)，其中 \(s\in S\)。我们将对 M 和 f 应用第 5 号的命题 5。~有 \(a(s,s')=\{s,s'\}\)，只要 \(s,s'\) 属于 S 且 \(m(s,s')\) 有限，因此存在一个映射 \(g:w\mapsto S_w\)，从 W 到 M，使得 \(g(w)=f(s_1)\cup\cdots\cup f(s_q)\)，即 \(S_w=\{s_1,\ldots,s_q\}\)，对于任意 \(w\in W\) 和任意约化分解 \((s_1,\ldots,s_q)\) 的 w 都如此。

推论 1. 对于任意子集 \(\mathbf{X}\)，它是 \(\mathbf{S}\) 的子集，\(\mathrm{W}_{\mathbf{X}}\) 作为 \(\mathbf{W}\) 的子群，由元素 \(w\)（属于 \(\mathbf{W}\)）组成，且这些元素满足 \(\mathrm{S}_w \subset \mathbf{X}\)。

若 \(w = s_1 \ldots s_q\)，其中 \(s_1, \ldots, s_q\) 属于 \(S\)，则 \(w^{-1} = s_q \ldots s_1\)；因此

\[
\mathrm{S} _ {w ^ {- 1}} = \mathrm{S} _ {w}.\tag{23}
\]

第 5 号的命题 4 表明，对于 \(S_{sw'} \subset \{s\} \cup S_{w'}\)，其中 \(s \in S\) 且 \(w' \in W\)，从而得到公式

\[
\mathrm{S} _ {w w ^ {\prime}} \subset \mathrm{S} _ {w} \cup \mathrm{S} _ {w ^ {\prime}}\tag{24}
\]

对 w 的长度作归纳即可证明。由 (23) 和 (24)，集合 U 由满足 \(w \in W\) 且 \(S_{w} \subset X\) 的元素组成，是 W 的一个子群；我们有 \(X \subset U \subset W_{X}\)，因此 \(U = W_{X}\)。

推论 2. 对于 S 的任意子集 X，都有 \(W_{X} \cap S = X\)。

这由推论 1 以及公式 \(S_s = \{s\}\)（其中 \(s\) 属于 \(S\)）得出。

推论 3. 集合 S 是 W 的极小生成集。

若 \(X \subset S\) 生成 \(W\)，则 \(W = W_X\)，因而由推论 2 有 \(X = S \cap W_X = S\)。

推论 4. 对于 S 的任意子集 X 以及任意 w 属于 \(W_{X}\)，w 关于 \(W_{X}\) 的生成集 X 的长度等于 \(l_{S}(w)\)。

设 \((s_{1},\ldots,s_{q})\) 是把 w 视为 W 的元素时的一个约化分解。我们有 \(w=s_{1}\ldots s_{q}\)，并且 \(s_{j}\in X\) 对 \(1\leqslant j\leqslant q\) 成立；此外，w 不可能是 \(q^{\prime}<q\) 个属于 \(X\subset S\) 的元素的乘积，这是由 \(q=l_{S}(w)\) 的定义所致。

定理 2. (i) 对于任意子集 \(X\)，它是 \(S\) 的子集，\((W_X, X)\) 是一个 Coxeter 系统。

\begin{enumerate}
\def\labelenumi{(\roman{enumi})}
\setcounter{enumi}{1}
\item
  设 \((\mathrm{X}_i)_{i\in \mathbf{I}}\) 是 S 的子集族。若 \(\mathrm{X} = \bigcap_{i\in \mathbf{I}}\mathrm{X}_i\)，则 \(\mathrm{W_X} = \bigcap_{i\in \mathbf{I}}\mathrm{W}_{\mathrm{X}_i}\)。
\item
  设 \(\mathbf{X}\) 和 \(\mathbf{X}'\) 是 \(\mathbf{S}\) 的两个子集。则 \(\mathrm{W_X} \subset \mathrm{W_{X'}}\)（分别地，\(\mathrm{W_X} = \mathrm{W_{X'}}\)），当且仅当 \(\mathbf{X} \subset \mathbf{X'}\)（分别地，\(\mathbf{X} = \mathbf{X'}\)）。
\end{enumerate}

X 的每个元素的阶都是 2，而 X 生成 \(W_{X}\)。设 \(x \in X\) 且 \(w \in W_{X}\)，满足 \(l_{\mathrm{X}}(xw) \leqslant l_{\mathrm{X}}(w) = q\)。由命题 7 的推论 4，有

\[
l _ {S} (x w) \leqslant l _ {S} (w) = q.
\]

设 \(x_{1},\ldots ,x_{q}\) 是 X 中满足 \(w = x_{1}\dots x_{q}\) 的元素。由于 (W,S) 满足交换条件（第 6 号的定理 1），存在整数 \(j\)，满足 \(1\leqslant j\leqslant q\)，且 \(xx_{1}\ldots x_{j - 1} = x_{1}\ldots x_{j - 1}x_{j}\)。因此，\((\mathrm{W_X},\mathrm{X})\) 满足交换条件，故是一个 Coxeter 系统（第 6 号的定理 1）。这就证明了 (i)。

断言 (ii) 和 (iii) 立即由命题 7 的推论 1 得出。

